\section{Reconstructing the law with every marginal retained}
\label{p13:transfer:section}

At a precision finer than $N^{-1}$, removing the exceptional windows is
too expensive. We retain their exact marginals instead. This also keeps
the repeated occurrence of a single Bernoulli variable inside the
reference law, where it can be evaluated exactly. The remaining task is
to control interactions between distinct windows and to sum the expansion
without applying a fixed-order estimate at an order growing with $N$.

We use the full admissible set $\mathcal I_N=\{x\in I_N:p_x>0\}$,
the exact marginals $p_x=\E J_x$, the centered coefficients
$c(I)=\E\prod_{x\in I}(J_x-p_x)$, and the relative nullity defined in
Section~\ref{p13:exact:section}. Throughout this section we abbreviate
$\rho_{\mathrm{rel}}(I)$ to $\rho(I)$; this convention uses the effective
local ranks and does not add the local rank defects again. Put
\[
 A_2=\sum_{|e|=2}|c(e)|,\qquad
 A_3=\sum_{\substack{|T|=3\\T\text{ connected}}}|c(T)|,
 \qquad p_* =\max_{x\in\mathcal I_N}p_x.
\]
The marginal and pair estimates give
$p_*\leq2p\ll_C N^{-1}$, $\mu=\sum_xp_x=O_C(1)$, and
$A_2\ll_{C,\varepsilon}N^{-2/3+\varepsilon}$.
The connection between the arithmetic estimates and the distribution is
the following bound:
\begin{equation}
 \|\Law(Z_N)-Q_{\mu_N,\alpha_N}\|_1
 \ll_{C,\varepsilon}N^\varepsilon A_3+N^{-4/3+\varepsilon}.
 \label{p13:transfer:announced-bound}
\end{equation}
We prove it in Theorem~\ref{p13:transfer:reconstruction}. The two terms
have different origins: the first records a genuinely connected
three-window interaction, while the second permits two pair interactions.

\paragraph{Three errors in the reconstruction.}
The retained class \(\mathscr D_2\) consists of direct families and
families \(e\cup U\) with a dependent pair \(e\) for which
\(S_e\) and the spaces \((S_z)_{z\in U}\) form a direct sum.
The dependent pair is then unique. Write
\(T_{\mathscr D_2}=\sum_{I\in\mathscr D_2}q(I)\Delta^{*|I|}\).
With the covariance-corrected Bernoulli measure \(B_2\) defined in
\eqref{p13:transfer:B2}, the proof separates
\[
 \begin{aligned}
 \|\Law(Z_N)-Q_{\mu_N,\alpha_N}\|_1
 &\le \|\Law(Z_N)-T_{\mathscr D_2}\|_1
       +\|T_{\mathscr D_2}-B_2\|_1\\
 &\qquad +\|B_2-Q_{\mu_N,\alpha_N}\|_1.
 \end{aligned}
\]
The first term is the stopping boundary, with the unexpanded suffix
kept in a positive conditional measure. The second consists of the
independent and pair deficits in
\eqref{p13:transfer:retained-difference}; the covariance signs are
retained before taking norms. In particular, an incompatible pair
is allowed in \(\mathscr D_2\): its joint probability is zero,
but its covariance is still present in \(B_2\).
The last term costs \(O(N^{-5/3+\varepsilon})\) by
Lemma~\ref{p13:transfer:conversion}; that rate concerns the full
measure \(B_2\), including every covariance.
Theorem~\ref{p13:transfer:finite} assembles the three estimates,
using fixed orders and the separately uniform tail of
Proposition~\ref{p13:transfer:uniform-tail} for the boundary.

\subsection{The exact Bernoulli reference and its conversion}

For any finite family of Bernoulli variables, without an independence
assumption, define $b_i=\delta_0+p_i\Delta$ and
\begin{equation}
 B_2=\mathop{*}_{i\in\mathcal I_N}b_i
       +\sum_{|e|=2}c(e)\Delta^{*2}*
                     \mathop{*}_{i\notin e}b_i.
 \label{p13:transfer:B2}
\end{equation}
This is a signed measure. Its first term is the law of a sum of
independent Bernoulli variables with the exact individual marginals;
its second term restores every covariance once.

\begin{lemma}[Conversion with heterogeneous marginals]
\label{p13:transfer:conversion}
Let $s_j=\sum_i p_i^j$, $C_2=\sum_{|e|=2}c(e)$, and
$\alpha=2C_2-s_2$. If $\mu\leq M$ and $p_*\leq1/4$, then
\begin{equation}
 \|B_2-Q_{\mu,\alpha}\|_1
 \leq C_M\{s_3+s_2^2+A_2(s_2+p_*)\}.
 \label{p13:transfer:conversion-bound}
\end{equation}
For the present model this is $O_{C,\varepsilon}(N^{-5/3+\varepsilon})$.
\end{lemma}

\begin{proof}
In the convolution algebra of finite signed measures, the series for
$\log_*(\delta_0+p_i\Delta)$ converges absolutely, since
$\|p_i\Delta\|_1=2p_i\leq1/2$. Summing it gives
\[
 \log_*\left(\mathop{*}_ib_i\right)
   =\mu\Delta-\frac{s_2}{2}\Delta^{*2}+E,
 \qquad \|E\|_1\leq\sum_{j\geq3}\frac{2^js_j}{j}
                   \leq\frac{16}{3}s_3.
\]
Exponentiation and $\exp_*(\mu\Delta)=\Pi_\mu$ imply
\[
 \left\|\mathop{*}_ib_i-\Pi_\mu
              +\frac{s_2}{2}\Pi_\mu*\Delta^{*2}\right\|_1
 \leq C_M(s_3+s_2^2).
\]
The product omitting a pair $e=\{i,j\}$ differs from
$\Pi_{\mu-p_i-p_j}$ by $O_M(s_2)$ in norm, and
\[
 \|\Pi_{\mu-p_i-p_j}-\Pi_\mu\|_1\leq2(p_i+p_j).
\]
Replacing each such product in \eqref{p13:transfer:B2} by $\Pi_\mu$
therefore costs $O_M(A_2(s_2+p_*))$. The exact identity
$\Var Z-\E Z=2C_2-s_2$ proves the lemma.
Here $s_2\leq p_*\mu=O_C(N^{-1})$ and
$s_3\leq p_*^2\mu=O_C(N^{-2})$, giving the stated application.
\end{proof}

\subsection{Two dense targets above a prescribed space}

The independent counts do not control a large rank gain: a family of
nullity $d$ carries an additional factor $2^d$. We now isolate the
conditional geometric estimate that pays for this factor. Only two
centers remain free. Every prime circuit meeting at least three
physical blocks must then meet one of the fixed sources, which makes
its coefficient enumerable by a divisor bound.

\begin{lemma}[Conditional dense pair estimate]
\label{p13:transfer:dense-targets}
Fix a number $s$ of source windows and an effective subspace $W$ of
their value space. Suppose $S_y\cap W=S_z\cap W=0$, and write
\[
 d=\dim W+\dim S_y+\dim S_z-\dim(W+S_y+S_z).
\]
For every fixed $\eta,\varepsilon>0$ there exists a fixed $\theta>0$
such that, uniformly in the sources and $W$,
\begin{equation}
 \sum_{\substack{y\ne z,\ y,z\text{ outside the sources}\\
             d\geq\eta L\\
             (y,z)\text{ nonaligned or }\dim(S_y\cap S_z)\leq\theta L}}
         p_yp_z2^d\ll_{s,C,\eta,\varepsilon}N^{-1+\varepsilon}.
 \label{p13:transfer:dense-target-bound}
\end{equation}
In particular the hypothesis on the pair is satisfied when $S_y,S_z$
are direct. For a primitive family of fixed size $r\geq3$, with a
prescribed collection of $r-2$ sources having at most $N^{a+\varepsilon}$
possibilities, its dense homogeneous mass is at most
\begin{equation}
 \sum_{\rho\geq\eta L}\Bigl(\prod_{i\in I}p_i\Bigr)2^{\rho(I)}
          \ll N^{-(r-1)+a+\varepsilon}.
 \label{p13:transfer:restricted-dense}
\end{equation}
\end{lemma}

The kernel defining $d$ projects injectively onto each target, so its
coordinate labels force many binary components in each target separately.
After regularization and removal of the weak aligned components, let
$m_t$ count the remaining components meeting target $t$, and let $n_{1,t}$
count those using exactly one of its $B=L+1$ coordinates. Then
$n_{1,t}\geq2m_t-B$. Above $d=(1/2+\delta)L$, this forces a positive
proportion of singleton occurrences in both targets, allowing the stronger
Pell or degree-one curve count. The resulting bounds are
$N^{1/2+\varepsilon'}$ target pairs in the lower stratum and
$N^{\varepsilon'}$ in the upper stratum; multiplying by $p_yp_z2^d$
gives $N^{-1+\varepsilon}$ after choosing the internal margins.
Section~\ref{p14:comp:transfer-dense-targets} supplies the separate
component counts, the coefficient enumeration from fixed sources, and
the removal of small aligned curves, including effective local sections.


\begin{lemma}[A dense compressed target]
\label{p13:transfer:compressed-target}
For fixed sources and a fixed channel of good pairs, the number of
compressed centers $X$ for which
$\dim(S(U_X)\cap W)\geq\eta L$ is $O_{s,C,\eta,\varepsilon}(N^\varepsilon)$.
\end{lemma}

\begin{proof}
The compression has fixed width $O(B)$ and an injective even-selection
map. In Proposition~\ref{p1:fixed:targets}, take $H$ larger than this
width. Its regularization loses a fixed small fraction of $\eta L$.
A prime above $H$ has at most one target occurrence, so all other
occurrences in its circuit are source occurrences. The circuit products
are divisors of fixed source values, and the fixed-source CRT and
split-product steps in that proposition are unchanged. The proof uses
the width bound and the target projection rank, not a target space of
dimension exactly $L$. The finitely many channel dilations and
$B^{O(1)}$ offset patterns are absorbed in the margins.
\end{proof}

\subsection{The exact weights of primitive families}

We use a small-kernel count to make a dense primitive rare before applying
the conditional estimate. For $Y\geq H$, put
\[
 K_Y(n)=\prod_{\substack{\ell>Y\\v_\ell(n)\ \mathrm{odd}}}\ell,\qquad
 a_Y(n)=\prod_{\ell>Y}\ell^{\lfloor v_\ell(n)/2\rfloor}.
\]
Then $n=s_Y(n)a_Y(n)^2K_Y(n)$, where $s_Y(n)$ is $Y$-smooth.

\begin{lemma}[Several small kernels in one window]
\label{p13:transfer:small-kernels}
Fix $A>2$, $0<b,u<1$, and an integer $m$ with $m(1-b)>1$.
For $Y=(\log N)^A$, the number of windows containing $m$ distinct
values with $K_Y(n)\leq N^b$ and $a_Y(n)^2\leq N^u$ is
\begin{equation}
 \ll N^{\max(1-1/A,(1+u)/2)+\varepsilon}.
 \label{p13:transfer:small-kernels-bound}
\end{equation}
Also, the set $\mathcal X_q$ of windows containing distinct $n_1,n_2$
with $K_H(n_1)K_H(n_2)\leq N^q$ has cardinality
$O(N^{q+\varepsilon})$, for fixed $q\geq0$.
\end{lemma}

\noindent The complete proof is given in Section~\ref{p14:comp:transfer-small-kernels}.


\begin{proposition}[Primitive mass at every fixed size]
\label{p13:transfer:primitive}
For every fixed integer $r\ge3$ and every $\varepsilon>0$,
\begin{equation}
 \sum_{\substack{I\subseteq\mathcal I_N,\ |I|=r\\I\ \mathrm{primitive}}}
         \Bigl(\prod_{x\in I}p_x\Bigr)2^{\rho(I)}
       \ll_{r,C,\varepsilon}N^{-r/2+\varepsilon}.
 \label{p13:transfer:primitive-bound}
\end{equation}
For each fixed $\eta>0$, primitive families with $\rho(I)\ge\eta L$
have $O_{r,C,\eta,\varepsilon}(N^\varepsilon)$ hosts and total
weight $O_{r,C,\eta,\varepsilon}(N^{1-r+\varepsilon})$.
These statements hold with exact heterogeneous marginals, including
irregular windows. The thresholds are chosen after the fixed parameters
and before $N,L$ and any deterministic subfamily being summed.
\end{proposition}

\begin{proof}
We first explain the dense estimate. The fixed-arity anchored argument
in Section~\ref{p14:comp:confinement} lifts effective relations to
physical even selections, regularizes their coordinate labels with
Lemma~\ref{p1:fixed:recovery}, and groups large primes by minimal
supports through one designated window. With $r$ fixed, each anchor
position belongs to at most a fixed number $D=D(r,\eta)$ of supports.
Primitivity excludes two exact units between any pair of windows.
Polynomial confinement therefore leaves only $o(B)$ groups larger
than $N^\tau$ for any fixed $\tau>0$. Erasing their anchor positions
costs $o(B)$ dimensions and leaves two values whose $B$-kernels have
product at most $N^{2D\tau}$.

Apply this construction separately at every window of the family.
All $r$ centres lie in the same deterministic set
$\mathcal X_{2D\tau}$, of size
$O(N^{2D\tau+\varepsilon'})$ by
Lemma~\ref{p13:half:small-kernels}.
Choose $\tau,\varepsilon'>0$ after $r,\eta,\varepsilon$ so that
$2rD\tau+r\varepsilon'<\varepsilon$.
This proves the asserted host count. The companion gives the
minimal-support construction, its incidence bound and the erasure
argument in detail; it does not assume the final host estimate.

The relation-space projection into each primitive block is injective,
so $\rho(I)\le\min_{x\in I}d_x\le L$. Since $p_x\le2p\ll_C N^{-1}$
and $2^L\ll_C N$, every weight is $O_{r,C}(N^{1-r})$.
Multiplying by the dense host count proves the second mass estimate.
For the complete sum, choose a new fixed $\xi>0$ within the final
margin and split at $\rho(I)=\xi L$.
Below that threshold, the independent-mass assertion of
Corollary~\ref{p13:half:connected} gives
$O(N^{-r/2+\xi+\varepsilon'})$.
Above it use $O(N^{1-r+\varepsilon'})$, with density parameter $\xi$.
As $r\ge3$, $1-r<-r/2$. Choose $\xi,\varepsilon'$ sufficiently small
to obtain \eqref{p13:transfer:primitive-bound}.

For a primitive family,
$c(I)=q(I)-\prod_{x\in I}p_x$ and
$q(I)$ is either zero or $2^{\rho(I)}\prod_{x\in I}p_x$.
Thus $q(I)+|c(I)|$ is at most twice the same weight, so the
corresponding actual joint and absolute centered sums have the same
powers. All orders and density parameters in this argument are fixed;
the all-cardinality estimates below still require their separate
certificate budgets.
\end{proof}

There is also an exact finite expansion:
\begin{equation}
 \Law\left(\sum_iJ_i\right)
   =\sum_{I\subseteq\mathcal I_N}c(I)\Delta^{*|I|}*
                                \mathop{*}_{i\notin I}b_i.
 \label{p13:transfer:exact-expansion}
\end{equation}
Indeed $\delta_{J_i}=b_i+(J_i-p_i)\Delta$; multiply these identities
by convolution and take expectations. The terms of cardinality one
vanish. The approximation $B_2$ retains exactly cardinalities zero and
two. We will not bound the absolute value of every remaining term
globally: the stopping identity retains a positive suffix and reduces
the required rank range.

\subsection{An independent extension estimate at all cardinalities}

For a fixed number $s$ of source windows and an effective subspace $W$
of their value space, let
\[
 E_W(u)=\sum_{\substack{U\subset\mathcal I_N\setminus\{\text{sources}\}\\
                      W,(S_x)_{x\in U}\text{ not direct}}}
                 u^{|U|}\prod_{x\in U}p_x,
 \qquad u>0.
\]
The weights here are products of marginal probabilities. They are not
conditional probabilities of dependent extensions.

\begin{lemma}[Uniform independent deficit]
\label{p13:transfer:independent-deficit}
For fixed $s,u,C$ and every $\varepsilon>0$,
\begin{equation}
 E_W(u)\ll_{s,u,C,\varepsilon}N^{-1/2+\varepsilon},
 \label{p13:transfer:independent-bound}
\end{equation}
uniformly in the sources and $W$. The sum includes all cardinalities.
\end{lemma}

\begin{proof}
Set $Y=B^2$ and
\[
 a(n)=\prod_{\substack{\ell>Y\\\ell\mid n}}(1+\ell^{-1/2}),\qquad
 A_x=\prod_{\substack{\ell>Y\\\ell\mid n\text{ for some }n\in V_x}}
                                    (1+\ell^{-1/2}).
\]
H\"older applied to the $B$ positions in a window gives
\[
 \sum_{N\leq x<2N}A_x
 \leq\prod_{i=-1}^{L-1}
          \left(\sum_{N\leq x<2N}a(x+i)^B\right)^{1/B}.
\]
Write $a(n)^B=\sum_{d\mid n}h_B(d)$, where $h_B$ is multiplicative,
supported on squarefree integers with primes exceeding $Y$, and
$h_B(\ell)=(1+\ell^{-1/2})^B-1\ll B\ell^{-1/2}$.
The coefficients are nonnegative, so
\[
 \sum_{n\leq3N}a(n)^B
 \leq3N\prod_{Y<\ell\leq3N}\left(1+\frac{h_B(\ell)}\ell\right)
 \leq3N\exp\left(O\left(B\sum_{m>Y}m^{-3/2}\right)\right)
 \ll N.
\]
Thus $\sum_xA_x\ll N$, with a constant independent of $B$.
The small-prime contribution satisfies
\[
 A_{\leq Y}=\prod_{\ell\leq Y}(1+\ell^{-1/2})
       \leq\exp(O(B/\log B))=N^{o(1)}.
\]
Chebyshev's bound and partial summation suffice for this estimate.

Let $P_0$ contain the small primes and the prime factors of the fixed
sources. Their total logarithmic size is $O_s(B\log N)$.
Using the same $Y=B^2$, we have
\[
 \sum_{\substack{\ell\in P_0\\\ell>Y}}\ell^{-1/2}
 \leq\frac{\sum_{\ell\in P_0}\log\ell}{\sqrt Y\log Y}
       =O_s(B/\log B).
\]
The primes up to $Y$ have the same bound by partial summation.
Consequently $A_0=\prod_{\ell\in P_0}(1+\ell^{-1/2})
=\exp(O_s(B/\log B))=N^{o(1)}$.
If $P_k$ also includes the primes of an
arbitrary list $x_1,\ldots,x_k$, then
\[
 \prod_{\ell\in P_k}(1+\ell^{-1/2})
       \leq A_0A_{\leq Y}\prod_{j=1}^k A_{x_j}.
\]
The elementary proof of the exceptional-target estimate
in Proposition~\ref{p1:arith:fixed-branching}, before restricting the
number of sources, gives
\begin{equation}
 D(x_1,\ldots,x_k)
 \ll B(s+k)+B\sqrt N\prod_{\ell\in P_k}(1+\ell^{-1/2}).
 \label{p13:transfer:explicit-exception}
\end{equation}
Contacts account for the first term. Outside them, a selected target
value whose odd prime factors all belong to $P_k$ has the form $a v^2$,
with $a$ squarefree and supported on $P_k$. Summing
$\sqrt{N/a}$ over these $a$ gives the second term, with $B$ possible
placements. This upper bound also includes local defects.

Put $w_x=up_x$, $w_*\ll N^{-1}$,
$M_0=\sum_xw_x=O_{u,C}(1)$, and
$M_1=\sum_xw_xA_x=O_{u,C}(1)$.
Every non-direct extension of size $h$ has an exceptional exposure
in each of its $h!$ orders. Sum over the possible first exceptional
position, allowing repeated centers for an upper bound. The first term
of \eqref{p13:transfer:explicit-exception} contributes at most
\[
 O(Bw_*)\sum_{h\geq1}\frac1{h!}
              \sum_{i=1}^h(s+i-1)M_0^{h-1}=O(B/N).
\]
The second contributes at most
\[
 O(B\sqrt Nw_*A_0A_{\leq Y})
 \sum_{h\geq1}\frac1{h!}\sum_{i=1}^h M_1^{i-1}M_0^{h-i}
       =N^{-1/2+o(1)}.
\]
The double sum is bounded by $e^{\max(1,M_0,M_1)}$.
The varying number of sources has therefore been averaged explicitly;
no constant from a fixed-order theorem has been used at a growing order.
\end{proof}

\begin{lemma}[A fixed number of extensions of a pair]
\label{p13:transfer:fixed-deficit}
For every fixed $k\geq1$,
\begin{equation}
 \sum_{|e|=2}|c(e)|
 \sum_{\substack{U\cap e=\varnothing,\ |U|=k\\
                  S_e,(S_z)_{z\in U}\text{ not direct}}}
                  \prod_{z\in U}p_z
 \ll_{k,C,\varepsilon}N^{-3/2+\varepsilon},
 \qquad S_e=\sum_{i\in e}S_i.
 \label{p13:transfer:fixed-pair-bound}
\end{equation}
\end{lemma}

\begin{proof}
If the new windows are not direct among themselves, choose a primitive
subfamily. A primitive of fixed size $h\geq2$ has independent mass
$O(N^{-h/2+\eta})$ by Theorem~\ref{p13:half:conditional} with no
prescribed sources. The largest such cost is $N^{-1+\eta}$.
After multiplication by $A_2$ and summation of the remaining independent
windows this is $N^{-5/3+O(\eta)}$.

Otherwise choose a minimal nonempty $A\subseteq U$ with
$S_e\cap\sum_{a\in A}S_a\ne0$.
If $|A|=1$, the total cost is the one-target deficit
of Proposition~\ref{p13:activity:deficit}.
If $h=|A|\geq2$, a relation meeting $S_e$ has a nonzero projection in
each of the $h$ target spaces, by minimality. Conditional on $e$,
Theorem~\ref{p13:half:conditional} gives $N^{h/2+\eta}$ possibilities.
The product of the $h$ target marginals costs $O_h(N^{-h})$; summing
the pair weights therefore gives $N^{-2/3-h/2+O(\eta)}$.
The remaining $k-h$ independent weights have bounded sum. There are
only finitely many choices of $h$ and roles because $k$ is fixed.
\end{proof}

\subsection{Uniform certificate tails and the deficit of all pairs}

We record the part of the prime-certificate argument that is uniform in
the number of blocks. A block may be an original window or one of a
fixed finite collection of geometric compressions of width $O(B)$.
Its center ranges over $O(N)$ values and its independent weight is
bounded by a fixed multiple of $p$. A relation of full support selects
a nonzero even vector in every block.

\begin{lemma}[Independent full-support tails]
\label{p13:transfer:independent-tail}
For every fixed weight $v>0$ and $a<2$, there is a fixed integer $s$
such that
\begin{equation}
 \sum_{\substack{|F|\geq s\\F\text{ has a full-support relation}}}
                  \prod_{x\in F}(vp_x)
 \ll N^{-a+\varepsilon}.
 \label{p13:transfer:tail-independent}
\end{equation}
The same conclusion holds with one block of any of the fixed compressed
types, weighted by $vp$. For this version the remaining factors are
$vp_x$.
\end{lemma}

\begin{proof}
Here are the parameter calculations in the certificate argument of
Lemmas~\ref{p1:core:moderate}--\ref{p1:core:generator}.
Choose $Y=(\log N)^{A_0}$ and a kernel threshold $T=N^b$, with
$A_0>1$ and $0<b<1$. A block containing a value whose odd kernel above
$Y$ is at most $T$ is marked. If
$1/2<\sigma<1$ and $A_0(1-\sigma)<1$, the Rankin count gives
$N^{1-\delta+o(1)}$ marked centers, where
$\delta=(1-b)(1-\sigma)>0$.
Every unmarked block in a full-support relation has two selected
values of kernel exceeding $T$. Their kernels are coprime, since
$Y$ exceeds the fixed block diameter.

For a certificate in which every receiving modulus is at most $O(NB)$,
the CRT and the exponential sum over cardinalities give independent
mass
\[
 N^{-2b+4/A_0+o(1)}.
\]
Here $m=\lceil2\log(O(N))/\log Y\rceil$ and
\[
 \log(B^m(m+1)!)=(4/A_0+o(1))\log N.
\]
If no such certificate exists, each unmarked block has a saturated
neighbor. Its directed degree is $N^{2/A_0+o(1)}$ in both directions,
by Lemma~\ref{p1:core:degrees}.
When $2/A_0+\delta<1$, the component generator of
Lemma~\ref{p1:core:generator}, tilted by $N^{\delta/2}$, stays bounded.
The mass at cardinality at least $s$ is consequently
$N^{-s\delta/2+o(1)}$.
Choose $A_0$ large and then $b$ close enough to one that
$2b-4/A_0>a$; choose an allowed $\sigma$ and then $s$ with
$s\delta/2>a$. All these choices are made before $N$.

The proof uses the actual prime divisors of selected values and a fixed
bound on the block width. Changing one block to a fixed geometric type
changes only constant factors and $B^{O(1)}$ offset choices. Its two
anchors, receiving congruences, and degree bounds obey the same estimates.
Thus the same parameter choices prove the compressed assertion.
\end{proof}

For $e=\{x,y\}$ put $E_e(u)=E_{S_e}(u)$ and
$D_2(u)=\sum_{|e|=2}|c(e)|E_e(u)$.

\begin{proposition}[The pair deficit over all cardinalities]
\label{p13:transfer:pair-deficit}
For fixed $u>0$ and every $\varepsilon>0$,
\begin{equation}
 D_2(u)\ll_{u,C,\varepsilon}N^{-3/2+\varepsilon}.
 \label{p13:transfer:pair-deficit-bound}
\end{equation}
In particular it is smaller than the $N^{-4/3}$ scale required below.
\end{proposition}

\begin{proof}
First consider extensions with an internal primitive family of new
windows. At each fixed size their independent mass is $O(N^{-1+\eta})$
by the half-exponent count; at all larger sizes apply
Lemma~\ref{p13:transfer:independent-tail} with any $a>1$.
Summing free complements costs at most $e^{u\mu}$. Their contribution
to $D_2$ is $O(A_2N^{-1+\eta})=O(N^{-5/3+2\eta})$.

For a direct extension, choose its minimal subfamily $A$ crossing $S_e$.
Use Lemma~\ref{p13:activity:pair-partition}. Weak and exceptional pairs
have total activity $O(N^{-1+\eta})$ and $O(N^{-1+\eta})$,
respectively; Lemma~\ref{p13:transfer:independent-deficit} bounds their
complete contribution by $N^{-3/2+O(\eta)}$.
There remain good pairs in a finite set of fixed channels. The cases
$|A|<s$ are covered by Lemma~\ref{p13:transfer:fixed-deficit}.
For $|A|\geq s$, compress the good pair geometrically. Minimality and
directness of $A$ give a full-support relation on the compressed block
and all blocks of $A$. The independent weight assigned to the
compressed center in Lemma~\ref{p13:transfer:independent-tail} is $p$,
whereas the required pair weight is at most $N^{-\kappa_b+o(1)}$,
$\kappa_b=2-1/b\geq5/3$. The ratio is at most
$N^{1-\kappa_b+o(1)}$. Choosing $a=1$ in that lemma bounds this tail
by $N^{-\kappa_b+o(1)}\leq N^{-5/3+o(1)}$.
Each good pair is represented by a fixed-width pattern and one center;
the $B^{O(1)}$ patterns are absorbed in the margins. Free complements
again cost a constant. Choose the finitely many internal margins after
$\varepsilon$ to conclude.
\end{proof}

\subsection{Every fixed connected order reduces to triples}

\begin{proposition}[Reduction of connected activities]
\label{p13:transfer:fixed-activities}
For every fixed $r\geq4$ and every $\varepsilon>0$,
\begin{equation}
 \sum_{\substack{|I|=r\\I\text{ connected}}}|c(I)|
 \ll_{r,C,\varepsilon}N^\varepsilon A_3+N^{-4/3+\varepsilon}.
 \label{p13:transfer:fixed-activity-bound}
\end{equation}
The incompatible part, for every fixed $r\geq4$, is
$O_{r,C,\varepsilon}(N^{-3/2+\varepsilon})$.
\end{proposition}

\begin{proof}
We first bound compatible families by their positive probabilities.
Let $T_3^+$ denote the sum of $q(T)$ over connected compatible triples.
For a primitive triple, $q(T)\leq2c(T)$. For a nonprimitive compatible
triple,
\[
 q(T)=c(T)+\sum_{e\subset T,\ |e|=2}c(e)p_{T\setminus e}
                         +\prod_{i\in T}p_i.
\]
The pair terms sum to the one-target deficit, and the independent term
has mass $O(N^{-3/2+\eta})$ by the half-exponent count. Therefore
\begin{equation}
 T_3^+\leq2A_3+O(N^{-3/2+\eta}).
 \label{p13:transfer:positive-triples}
\end{equation}
Families containing such a triple have mass $N^\eta T_3^+$ after
positive propagation~\eqref{p13:exact:propagation}.
Families containing a primitive of size at least four have mass
$N^{-4/3+\eta}$, by Proposition~\ref{p13:transfer:primitive} and the
same propagation. All seed sizes here are bounded by the fixed $r$.

Consider a remaining compatible connected family. It contains a
dependent pair $e$; otherwise a minimal dependent subfamily would be
one of the seeds just treated. Put $U=S_e$ and choose a nonempty
$A\subset I\setminus e$ minimal with
$U\cap\sum_{a\in A}S_a\ne0$. A single target gives a connected triple,
so $h=|A|\geq2$. If $A$ contains a proper primitive $P$, then
$U\cap\sum_{i\in P}S_i=0$ by minimality. Thus
$q(e\cup P)=q(e)q(P)$. Only $|P|=2$ remains, since larger primitives
were already treated. Compatible dependent pairs satisfy $q(e)\leq2c(e)$,
so the total seed mass is at most $4A_2^2$.
Propagate it to $I$.

Suppose next that $A$ is direct. The nullity $d$ of the compressed
family $U,(S_a)_{a\in A}$ is at most $L$, and
\[
 q(e\cup A)=q(e)\Bigl(\prod_{a\in A}p_a\Bigr)2^d.
\]
For $d<\eta L$, the conditional half-exponent count of the $h$ targets
gives a cost at most $A_2N^{-h/2+O(\eta)}$, hence
$N^{-5/3+O(\eta)}$. For $d\geq\eta L$, fix $h-2$ members of $A$.
The last two targets are each direct from the sum of $U$ and these
fixed targets, and are direct from one another. Otherwise there would
be an internal dependence in $A$ or an earlier crossing of $U$.
Lemma~\ref{p13:transfer:dense-targets} gives conditional weighted cost
$N^{-1+\eta}$. The independent weights of the $h-2$ additional sources
sum to a constant; multiplication by the pair mass gives
$N^{-5/3+O(\eta)}$ again. If $A$ is itself primitive of size at least
three, it was already one of the excluded seeds.

It remains that $A=f=\{y,z\}$ is a dependent pair. By minimality,
$U\cap S_y=U\cap S_z=0$. Put
\[
 d_e=\dim(S_x\cap S_{x'}),\quad d_f=\dim(S_y\cap S_z),\quad
 \delta=\dim(U\cap S_f),\qquad e=\{x,x'\}.
\]
If $\delta<\eta L$, then
$q(e\cup f)=q(e)q(f)2^\delta$, and the total is
$O(N^\eta A_2^2)$.
For $\delta\geq\eta L$, choose $\theta$ from the dense-target lemma.
If $f$ is nonaligned or $d_f\leq\theta L$, that lemma gives
$N^{-5/3+O(\eta)}$ after summing $e$.
If $f$ is a good pair of a fixed channel, fix $e$ and use
Lemma~\ref{p13:transfer:compressed-target}; there are only
$N^\eta$ choices for the compressed center of $f$. Linear algebra gives
\begin{equation}
 d_f+\delta\leq L,
 \qquad q(e\cup f)\leq p_*q(e).
 \label{p13:transfer:two-pair-rank}
\end{equation}
Indeed the projection of $U\cap(S_y+S_z)$ modulo $S_y$ is injective,
and its image has dimension at most $\dim S_z-d_f$.
This case also costs $N^{-5/3+O(\eta)}$.

Every aligned pair with rank greater than $\theta L$ has bounded
channel height, by Lemma~\ref{p1:arith:alignedrank}. Outside the good
pairs at those finitely many heights, the normalization estimates in
Theorem~\ref{p1:arith:pairs} and
 Lemma~\ref{p1:arith:fixed-channel-mass} give at most
$N^{1/2+\eta}$ exceptional aligned pairs. The endpoint pairs of height
two and the pairs meeting an irregular window satisfy the same bound;
height one has rank $o(L)$ when compatible.
Call this exceptional collection $\mathcal B$.
If only $f$ belongs to $\mathcal B$, exchange $e$ and $f$ and use one
of the preceding cases. This exchange is valid because a target in
$e$ intersecting $S_f$ would form a connected triple, already removed.
If both belong to $\mathcal B$, there are at most $N^{1+2\eta}$ hosts.
The symmetric forms of \eqref{p13:transfer:two-pair-rank} give
\[
 d_e+d_f+\delta\leq L+\min(d_e,d_f)\leq\tfrac32L+o(L),
\]
where the final inequality uses
Theorems~\ref{p13:rig:nonaligned} and~\ref{p13:rig:aligned}.
The joint probability is at most $N^{-5/2+o(1)}$, so this contribution
is $N^{-3/2+O(\eta)}$. These cases exhaust the possible minimal sets
$A$. Positive propagation adds the remaining windows, proving the
compatible part of \eqref{p13:transfer:fixed-activity-bound}.

For incompatible families we use induction on $r$, starting with the
triple estimate. The exact identity is
\begin{equation}
 c(I)=-\prod_{i\in I}p_i-
          \sum_{\substack{J\subsetneq I\\|J|\geq2}}
                     c(J)\prod_{i\in I\setminus J}p_i.
 \label{p13:transfer:incompatible-expansion}
\end{equation}
The independent term over connected $r$-tuples is
$O(N^{-r/2+\eta})$. The pair terms cost $O(N^{-3/2+\eta})$ by
Lemma~\ref{p13:transfer:fixed-deficit}.
For a fixed proper $J$ of size at least three, connectedness of $I$
requires a non-direct extension of its total space; its independent
cost is $N^{-1/2+\eta}$ by
Lemma~\ref{p13:transfer:independent-deficit}.
Canonical factorization of $c(J)$ has no singleton factors.
By induction a factor of size at least three costs
$N^\eta A_3+N^{-4/3+\eta}$; a decomposition consisting only of pairs
has at least two of them and costs $O(A_2^2)$.
There are finitely many partitions at each fixed size. Multiplying by
the independent extension cost proves
\eqref{p13:transfer:fixed-activity-bound} at order $r$.
Inserting the arithmetic bound $A_3\ll N^{-41/30+\eta}$ from
Corollary~\ref{p13:third:activity} into the same induction gives the
stated $N^{-3/2+\varepsilon}$ incompatible activity. No law
approximation or lower bound has been used in this induction.
\end{proof}

\subsection{A uniform tail after 256 blocks}

The preceding proposition has constants depending on its fixed order.
We now use the certificate generator to handle all larger orders in
the stopped expansion. The number $256$ is a sufficient fixed cutoff,
not an optimal threshold.

\begin{lemma}[Certificate budgets with several anchors]
\label{p13:transfer:anchor-budgets}
Choose an even integer $t\geq2$, $A_0>1$, $0<b<1$, and
$1/2<\sigma<1$ with $A_0(1-\sigma)<1$. Put
\[
 \delta=(1-b)(1-\sigma),\qquad \theta=t/A_0.
\]
Suppose each unmarked block in a family has a relation selecting at
least $t$ values in that block; the relation may depend on the block.
Mark a block when some value has odd kernel above
$Y=(\log N)^{A_0}$ at most $N^b$.
Under any fixed independent weight per block, the moderate certificates
have total mass
\begin{equation}
 N^{-tb+2t/A_0+o(1)}.
 \label{p13:transfer:moderate-budget}
\end{equation}
If $\theta+\delta<1$, the saturated families of size at least $s$
have total independent mass
\begin{equation}
 N^{-s\delta/2+o(1)}.
 \label{p13:transfer:saturated-budget}
\end{equation}
\end{lemma}

\begin{proof}
This is the calculation in Lemma~\ref{p13:transfer:independent-tail}
with $t$ anchors. Their pairwise coprime kernels have product exceeding
$N^{tb}$ and at most
$m_t=\lceil t\log(O(N))/\log Y\rceil$ prime factors.
The CRT bound and exponential cardinality sum have combinatorial cost
$B^{m_t}(m_t+1)!=N^{2t/A_0+o(1)}$, proving
\eqref{p13:transfer:moderate-budget}.
The saturated graph has degree $N^{\theta+o(1)}$ in both directions;
the number of marked vertices is $N^{1-\delta+o(1)}$.
After tilting by $N^{\delta/2}$, the three terms of the component
generator of Lemma~\ref{p1:core:generator} are
\[
 N^{-\delta/2+o(1)},\quad
 N^{-1+\theta+\delta/2+o(1)},\quad
 N^{-1+\theta+\delta+o(1)}.
\]
They remain bounded, and in fact tend to zero, under the stated
condition. Removing the tilt proves
\eqref{p13:transfer:saturated-budget}. The inequality $p_x\leq2p$
is absorbed by doubling the fixed weight. The arguments concern true
prime factors and permit contacts and irregular source windows.
\end{proof}

For a family $D$, let $\mathcal K(D)$ be the kernel of the sum map
$\bigoplus_{x\in D}S_x\to\sum_{x\in D}S_x$, and set
\[
 d_x(D)=\dim\operatorname{proj}_{S_x}\mathcal K(D)
       =\rho(D)-\rho(D\setminus\{x\}).
\]
Call $D$ a $10$-core when each $d_x(D)\geq9$.
The effective projection admits a linear section in the formal even
coordinates. Setting nine pivot coordinates to one gives a word of
weight at least nine; even parity makes its weight at least ten.
Thus the preceding budgets apply with $t=10$.
For a compatible family Fourier expansion gives exactly
\begin{equation}
 \frac{c(D)}{q(D)}
 =\P_{v\ \mathrm{uniform\ in}\ \mathcal K(D)}
       (\operatorname{proj}_{S_x}v\ne0\text{ for every }x\in D)
 \geq1-|D|2^{-9}.
 \label{p13:transfer:core-ratio}
\end{equation}
In particular $q(D)\leq2c(D)$ for $|D|<256$.

\begin{proposition}[Uniform connected and primitive tails]
\label{p13:transfer:uniform-tail}
For fixed $w>0$, every $\varepsilon>0$, and every $K_N=o(L)$,
\begin{align}
 \sum_{\substack{I\text{ connected},\ 256<|I|\leq K_N\\
                         \rho^*(I)\leq2L}}
              w^{|I|}|c(I)|&\ll N^{-4/3+\varepsilon},
                    \label{p13:transfer:connected-tail}\\
 \sum_{\substack{P\text{ primitive}\\|P|>256}}
              w^{|P|}|c(P)|&\ll N^{-4/3+\varepsilon}.
                    \label{p13:transfer:primitive-tail}
\end{align}
The second sum includes all cardinalities. Here
$\rho^*(I)=\max\{\rho(A):A\subseteq I,\ q(A)>0\}$.
\end{proposition}

\begin{proof}
For compatible $10$-cores of nullity at most $2L$, choose
\[
 t=10,\quad A_0=12,\quad b=31/60,\quad
 \sigma=23/25,\quad\delta=29/750,\quad\theta=5/6.
\]
The condition $\theta+\delta<1$ is strict. Multiplying the independent
budgets by $2^{2L}=O_C(N^2)$ gives, at sizes at least $256$,
\begin{equation}
 \sum_D w^{|D|}q(D)
 \ll N^{-3/2+o(1)}+N^{-1106/375+o(1)}.
 \label{p13:transfer:large-core}
\end{equation}
Indeed $2-10(31/60)+20/12=-3/2$ and
$2-256(29/1500)=-1106/375$.

If $\rho^*(I)\leq L/60$, the centered expansion gives
\[
 |c(I)|\leq2^{|I|}\Bigl(\prod_{x\in I}p_x\Bigr)2^{L/60}.
\]
A nonzero centered coefficient has a full-support Fourier relation,
so two anchors are available in every block. Use
\[
 t=2,\quad A_0=12,\quad b=17/20,\quad
 \sigma=23/25,\quad\delta=3/250,\quad\theta=1/6.
\]
Starting at cardinality $256$, the two costs after the rank factor are
\begin{equation}
 N^{-27/20+o(1)}+N^{-2279/1500+o(1)}.
 \label{p13:transfer:low-rank-tail}
\end{equation}
The four costs in \eqref{p13:transfer:large-core} and
\eqref{p13:transfer:low-rank-tail} leave the exact margins
\[
 \begin{aligned}
 \frac32-\frac43&=\frac16,&
 \frac{1106}{375}-\frac43&=\frac{202}{125},\\
 \frac{27}{20}-\frac43&=\frac1{60},&
 \frac{2279}{1500}-\frac43&=\frac{93}{500}.
 \end{aligned}
\]
Thus the smallest margin is $1/60$. For the fixed certificate parameters
and fixed weight $w$, every displayed $o(1)$ loss is less than $1/120$
for sufficiently large $N$; no certificate parameter is allowed to vary
with $N$.

For the remaining connected $I$, choose deterministically a compatible
$A\subseteq I$ attaining $\rho^*(I)>L/60$. Delete blocks of projection
dimension at most eight until a $10$-core $D$ remains. Since
$8K_N=o(L)$, its nullity exceeds $L/120$ for large $N$, so $|D|\geq2$.
With $r=|I|$ and $h=|D|$, the centered expansion and maximality of
$\rho^*(I)$ imply
\begin{equation}
 |c(I)|\leq2^r2^{8(r-h)}
               \Bigl(\prod_{x\in I\setminus D}p_x\Bigr)q(D).
 \label{p13:transfer:peeling}
\end{equation}
The elements outside $A$ are already in this product; there is no
enumeration of deletion orders. All fixed powers of two can be absorbed
in the independent weight of the complements.

If $h\geq256$, sum the complements freely and apply
\eqref{p13:transfer:large-core}. If $3\leq h<256$, then
$q(D)\leq2c(D)$ by \eqref{p13:transfer:core-ratio}.
Proposition~\ref{p13:transfer:fixed-activities}, the arithmetic triple
bound, and canonical factorization show that the sum of $|c(D)|$ at
each such fixed size is $O(N^{-1+\eta})$. Because $D$ is proper in
the connected $I$, its complement is a non-direct independent extension.
Lemma~\ref{p13:transfer:independent-deficit} adds $N^{-1/2+\eta}$.
If $h=2$, then $D=e$ is a compatible dependent pair and $q(e)\leq2c(e)$;
the complement is exactly the pair deficit of
Proposition~\ref{p13:transfer:pair-deficit}. These bounds prove
\eqref{p13:transfer:connected-tail}.

A primitive family has nullity at most $L$. If it is incompatible,
$|c(P)|=\prod_Pp_x$ and the independent two-anchor budgets suffice.
If it is compatible with nullity at most $L/60$, use
\eqref{p13:transfer:low-rank-tail}. Otherwise each projection has
dimension $\rho(P)>L/60$, so it is itself a $10$-core for large $N$.
Apply \eqref{p13:transfer:large-core} directly, without an upper
cardinality cutoff. This proves \eqref{p13:transfer:primitive-tail}.
\end{proof}

\subsection{The stopped expansion and its finite input}

Let $\mathscr D_2$ consist of the direct families and the families
$e\cup U$ for which $e$ is any dependent pair and $S_e,(S_z)_{z\in U}$
are direct. Affine incompatibility of $e$ is allowed. This class is
hereditary; every non-direct member contains exactly one dependent pair.
For a fixed ordering of $\mathcal I_N$, its boundary contains the first
excluded prefixes. A boundary family has relative nullity at most $2L$:
its retained predecessor has at most $L$, and adding one block adds
at most $L$. In particular every subfamily of a boundary family has
compatible nullity at most $2L$.

\begin{theorem}[Finite transfer to the complete law]
\label{p13:transfer:finite}
Fix $1<\beta\leq4/3$. Suppose that, for each of the fixed integers
$3\leq r\leq256$ and every $\varepsilon>0$,
\begin{equation}
 \sum_{\substack{|I|=r,\ I\text{ connected}\\\rho^*(I)\leq2L}}
                   |c(I)|\ll_{r,C,\varepsilon}N^{-\beta+\varepsilon}.
 \label{p13:transfer:finite-hypothesis}
\end{equation}
Then
\begin{equation}
 \|\Law(Z_N)-Q_{\mu_N,\alpha_N}\|_1
                  \ll_{C,\varepsilon}N^{-\beta+\varepsilon}.
 \label{p13:transfer:finite-conclusion}
\end{equation}
Only these $254$ fixed orders are assumed in
\eqref{p13:transfer:finite-hypothesis}; the larger orders are supplied
by Proposition~\ref{p13:transfer:uniform-tail}.
\end{theorem}

\begin{proof}
Fix a weight $w>0$. The hypothesis at sizes at most $256$ and
\eqref{p13:transfer:connected-tail} give, for $K_N=o(L)$,
\begin{equation}
 T:=\sum_{\substack{I\text{ connected},\ 3\leq|I|\leq K_N\\
                             \rho^*(I)\leq2L}}
                  w^{|I|}|c(I)|\ll N^{-\beta+\eta}.
 \label{p13:transfer:connected-range}
\end{equation}
The finitely many factors $w^r$ at $r\leq256$ only change constants.
For a primitive $P$, $\rho(P)\leq L$ and
$\prod_{x\in P}p_x\leq|c(P)|$, because
$c(P)=q(P)-\prod_Pp_x$ and $q(P)$ is either zero or at least twice
the independent product. Together with
\eqref{p13:transfer:primitive-tail}, this gives
\begin{equation}
 \sum_{\substack{P\text{ primitive}\\|P|\geq3}}
                     w^{|P|}|c(P)|\ll N^{-\beta+\eta}
 \label{p13:transfer:all-primitives}
\end{equation}
over all cardinalities.

Put $T_{\mathscr D_2}=\sum_{I\in\mathscr D_2}q(I)\Delta^{*|I|}$.
The finite expansion of $B_2$ yields the identity
\begin{equation}
 T_{\mathscr D_2}-B_2
 =-\sum_{I\notin\mathscr D_2}\Delta^{*|I|}
       \left[\prod_{i\in I}p_i+
          \sum_{\substack{e\subset I\\|e|=2}}
                   c(e)\prod_{i\in I\setminus e}p_i\right].
 \label{p13:transfer:retained-difference}
\end{equation}
For the independent term, choose a primitive subfamily of a discarded
family. If it has at least three members, use
\eqref{p13:transfer:all-primitives}; if it is a pair $e$, the complement
must be a non-direct extension of $S_e$, and
$\prod_{i\in e}p_i\leq|c(e)|$.
The pair terms in \eqref{p13:transfer:retained-difference} are bounded
by the same all-cardinality pair deficit. Free complements have total
weight at most $e^{2\mu}$. Consequently
\begin{equation}
 \|T_{\mathscr D_2}-B_2\|_1\ll N^{-\beta+\eta}.
 \label{p13:transfer:retained-error}
\end{equation}

Take $K_N=(5+o(1))\log N/\log\log N$.
The exact stopping identity from Section~\ref{p1:stop:section} gives
\[
 \|\Law(Z_N)-T_{\mathscr D_2}\|_1
       \leq\sum_{I\in\partial\mathscr D_2}2^{|I|}q(I).
\]
For $|I|\leq K_N$, expand $q(I)$ in centered coefficients. The terms
of size zero and two are the independent and pair deficits already
bounded. Every remaining nonzero centered term has either a connected
component of size at least three or at least two pair components.
Canonical factorization, \eqref{p13:transfer:connected-range}, and
the free independent complements bound their sum by a constant times
\[
 e^{4A_2+T}-1-4A_2=O(T+A_2^2).
\]
Here $T=o(1)$ and $A_2^2\ll N^{-4/3+\eta}$; the centered coefficient
of a singleton is zero. This is where the linear contribution of one
pair is subtracted before taking an absolute bound.

On the larger boundary, $q(I)\leq2^{2L}\prod_Ip_x$, so
\[
 \sum_{\substack{I\in\partial\mathscr D_2\\|I|>K_N}}2^{|I|}q(I)
 \ll_C N^2\sum_{r>K_N}\frac{(2\mu)^r}{r!}=N^{-3+o(1)}.
\]
Combining these bounds, \eqref{p13:transfer:retained-error}, and
Lemma~\ref{p13:transfer:conversion} proves
\eqref{p13:transfer:finite-conclusion}. All internal margins are chosen
after the final $\varepsilon$; none is made to vary with $N$.
\end{proof}

\begin{theorem}[Reconstruction from triple activity]
\label{p13:transfer:reconstruction}
With the exact mean and second factorial cumulant,
\begin{equation}
 \|\Law(Z_N)-Q_{\mu_N,\alpha_N}\|_1
       \ll_{C,\varepsilon}N^\varepsilon A_3+N^{-4/3+\varepsilon}.
 \label{p13:transfer:reconstruction-bound}
\end{equation}
In particular the arithmetic triple estimate gives the upper bound
$O_{C,\varepsilon}(N^{-4/3+\varepsilon})$ for the full model.
\end{theorem}

\begin{proof}
In the proof of Theorem~\ref{p13:transfer:finite}, replace the finite
activity input by Proposition~\ref{p13:transfer:fixed-activities}.
At sizes $3,\ldots,256$ its total cost is
$N^{\varepsilon'}A_3+N^{-4/3+\varepsilon'}$; the larger sizes have the
second cost by Proposition~\ref{p13:transfer:uniform-tail}.
The primitive deficits, pair deficits, and boundary expansion obey
the same bounds. In the component generator,
\[
 T\ll N^{\varepsilon'}A_3+N^{-4/3+\varepsilon'},\qquad
 e^{4A_2+T}-1-4A_2=O(T+A_2^2).
\]
The established arithmetic estimate makes $T=o(1)$, so this passage
loses no new fixed power of $N$. The large boundary costs $N^{-3+o(1)}$
and conversion costs $N^{-5/3+\varepsilon'}$.
To make the choices explicit, denote the requested final margin by
$\varepsilon_{\mathrm{out}}=\varepsilon$ and choose
$0<\tau<\min(\varepsilon_{\mathrm{out}}/4,1/960)$.
Invoke each completed fixed-order estimate with sufficiently small fixed
input margins that its total displayed loss is at most $\tau$.
For example, use the pair estimate with loss $\tau/2$ before squaring
$A_2$. Only the fixed orders $3,\ldots,256$ occur in this choice;
the larger orders use the uniform certificate bounds, not constants
from a fixed-order theorem. With $\varepsilon'=\tau$ the preceding
argument proves \eqref{p13:transfer:reconstruction-bound}, since
$\tau<\varepsilon_{\mathrm{out}}$.

Finally invoke Corollary~\ref{p13:third:activity} with loss $\tau$.
The triple term is then at most $N^{-41/30+2\tau}$, and
\[
 \frac{41}{30}-2\tau-\frac43=\frac1{30}-2\tau>0.
\]
The conversion and large-boundary exponents exceed $4/3$ by $1/3$
and $5/3$, respectively. These comparisons give the announced upper
bound. All accuracy choices, source counts and fixed activity orders used here
are fixed before $N$; the growing cardinality cutoff $K_N$ is controlled
separately.
\end{proof}

The method retains three separate forms of control. Exact Bernoulli
marginals handle repetitions of one event; arithmetic counts bound
interactions among a fixed number of distinct windows; the certificate
generator and the positive stopped suffix sum the growing part of the
expansion. Neither bounded fixed moments alone nor a finite list of
cumulants replaces the third step.
